% ══════════════════════════════════════════════════════════════
\section{CENA 1 --- Abertura\hfill\normalfont\small\color{corTempo}0:00 -- 2:00}
% ══════════════════════════════════════════════════════════════

\tempo{0:00 -- 0:15}

\begin{audiocue}
\faMusic\ Trilha: intro energética, batida eletrônica crescente. Fade out suave ao entrar a voz.
\end{audiocue}

\begin{visual}
\faFilm\ TAKE INICIAL: close em chip moderno com LEDs piscando. Zoom rápido para dentro:
transistores, trilhas de cobre, elétrons em movimento. Gráfico animado da Lei de Moore
ao fundo -- curva exponencial subindo.
\end{visual}

\tempo{0:15 -- 0:45}

\begin{fala}
Você sabia que o processador do seu computador pode ficar mais quente do que
a superfície de uma chapa de fogão?
\end{fala}

\begin{visual}
\faFilm\ Chip esquentando gradualmente: cores mudam de cinza para laranja e vermelho.
Termômetro ao lado: 100, 120, 140, 160 {\textdegree}C.
\end{visual}

\begin{fala}
A Lei de Moore -- aquela regra que diz que os chips dobram de capacidade a cada
dois anos -- tem um preço que pouca gente comenta: \textbf{densidade de potência
absurda}. E densidade de potência significa \textbf{calor}.
\end{fala}

\tempo{0:45 -- 1:20}

\begin{visual}
\faFilm\ Split-screen: à esquerda, chip clássico com hotspot vermelho brilhando.
À direita, esfera translúcida (Esfera de Bloch) tremendo suavemente -- o qubit perturbado.
\end{visual}

\begin{fala}
Mas essa batalha contra o calor não é só coisa de computador tradicional.
Na computação quântica -- que promete resolver problemas impossíveis para qualquer
máquina atual -- o calor é um inimigo ainda mais traiçoeiro.
Ele não derrete nada. Mas comete erros no silêncio.
\end{fala}

\begin{tela}
\faDesktop\ \textbf{Texto na tela:}
\begin{center}
``O calor não avisa. Ele age.''
\end{center}
\end{tela}

\tempo{1:20 -- 2:00}

\begin{fala}
Neste vídeo, a gente vai usar uma ferramenta do cálculo -- as
\textbf{integrais múltiplas} -- para medir, prever e entender esses efeitos
nos dois mundos: no chip de silício e no qubit quântico.
E no final, a gente conecta os dois. Porque o inimigo é o mesmo.
\end{fala}

\begin{tela}
\faDesktop\ \textbf{Texto aparece letra por letra:}
\begin{center}
\large\itshape Integrais Duplas e Triplas:\\
a matemática que luta contra o calor.
\end{center}
\end{tela}

\begin{audiocue}
\faMusic\ Trilha sobe brevemente e cai para fundo suave e contínuo.
\end{audiocue}